% ═══════════════════════════════════════════════════════════════
% SLIDES: El tiempo y los lugares
% Spanisch Klasse 8 - Unidad 6, DS 1 (90 Min)
% ═══════════════════════════════════════════════════════════════

\documentclass[aspectratio=169,11pt]{beamer}

\usepackage[utf8]{inputenc}
\usepackage[T1]{fontenc}
\usepackage[ngerman]{babel}
\usepackage{helvet}
\usepackage{tikz}
\usepackage{tabularx}
\usepackage{booktabs}

% ═══════════════════════════════════════════════════════════════
% THEME & FARBEN
% ═══════════════════════════════════════════════════════════════

\usetheme{default}
\usecolortheme{default}
\setbeamertemplate{navigation symbols}{}
\setbeamertemplate{footline}[frame number]

\definecolor{spanisch}{HTML}{C41E3A}
\definecolor{hellgrau}{HTML}{F5F5F5}
\definecolor{dunkelgrau}{HTML}{333333}

\setbeamercolor{frametitle}{fg=white,bg=spanisch}
\setbeamercolor{title}{fg=white,bg=spanisch}
\setbeamercolor{structure}{fg=spanisch}
\setbeamercolor{block title}{fg=white,bg=spanisch}
\setbeamercolor{block body}{bg=hellgrau}

\setbeamerfont{frametitle}{series=\bfseries}
\setbeamerfont{title}{series=\bfseries}

% Befehle
\newcommand{\es}[1]{\textcolor{spanisch}{\textbf{#1}}}
\newcommand{\de}[1]{\textcolor{dunkelgrau}{\textit{#1}}}
\newcommand{\phase}[1]{\begin{center}\colorbox{spanisch}{\textcolor{white}{\large\bfseries #1}}\end{center}}

% ═══════════════════════════════════════════════════════════════
% DOKUMENT
% ═══════════════════════════════════════════════════════════════

\begin{document}

% ---------------------------------------------------------------
% TITEL
% ---------------------------------------------------------------
\begin{frame}
\begin{center}
\vspace{1cm}
{\Huge\bfseries\textcolor{spanisch}{El tiempo y los lugares}}\\[0.5cm]
{\Large Das Wetter und die Orte}\\[1cm]
{\normalsize Spanisch Klasse 8 · Unidad 6 · Qué pasa 1, S. 102-103}
\end{center}
\end{frame}

% ---------------------------------------------------------------
% LERNZIELE
% ---------------------------------------------------------------
\begin{frame}{Hoy aprendemos...}
\begin{itemize}
\setlength\itemsep{1em}
\item[\textcolor{spanisch}{\textbullet}] \textbf{Wetter beschreiben:} ¿Qué tiempo hace?
\item[\textcolor{spanisch}{\textbullet}] \textbf{Orte im Viertel benennen:} el parque, la biblioteca...
\item[\textcolor{spanisch}{\textbullet}] \textbf{Aktivitäten mit Wetter verbinden:} Cuando hace sol, voy al parque.
\item[\textcolor{spanisch}{\textbullet}] \textbf{Kleidung und Wetter:} ¿Qué ropa llevas cuando hace frío?
\end{itemize}
\end{frame}

% ═══════════════════════════════════════════════════════════════
% PHASE 1: EINSTIEG (10 Min)
% ═══════════════════════════════════════════════════════════════

\begin{frame}{}
\phase{Einstieg: ¿Qué veis?}
\vspace{0.5cm}
\begin{center}
{\Large Öffnet das Buch auf Seite 102-103.}\\[1cm]
{\LARGE\bfseries 6 Fotos – ¿Qué veis?}\\[0.5cm]
{\normalsize Beschreibt frei, was ihr seht (Deutsch oder Spanisch).}
\end{center}
\end{frame}

% ═══════════════════════════════════════════════════════════════
% PHASE 2: INPUT WETTER (15 Min)
% ═══════════════════════════════════════════════════════════════

\begin{frame}{}
\phase{Input: El tiempo}
\end{frame}

\begin{frame}{¿Qué tiempo hace? \de{– Wie ist das Wetter?}}
\begin{columns}
\column{0.5\textwidth}
\textbf{Mit \es{hace}:}
\begin{itemize}
\item \es{Hace sol.} \de{Es ist sonnig.}
\item \es{Hace calor.} \de{Es ist heiß.}
\item \es{Hace frío.} \de{Es ist kalt.}
\item \es{Hace viento.} \de{Es ist windig.}
\end{itemize}

\column{0.5\textwidth}
\textbf{Mit \es{hace} + Bewertung:}
\begin{itemize}
\item \es{Hace buen tiempo.} \de{Gutes Wetter.}
\item \es{Hace mal tiempo.} \de{Schlechtes Wetter.}
\item \es{Hace 20 grados.} \de{Es sind 20 Grad.}
\end{itemize}
\end{columns}
\end{frame}

\begin{frame}{Wetter mit eigenen Verben}
\begin{center}
\begin{tabular}{ll}
\es{Llueve.} & \de{Es regnet.} \\[0.8em]
\es{Nieva.} & \de{Es schneit.} \\[0.8em]
\es{Hay nubes.} & \de{Es ist bewölkt.} \\
\end{tabular}
\end{center}

\vspace{1cm}

\begin{block}{Merke}
\es{Hace} + Substantiv (sol, frío, calor...)\\
Aber: \es{Llueve} / \es{Nieva} sind eigenständige Verben!
\end{block}
\end{frame}

\begin{frame}{Wiederholen wir!}
\begin{center}
{\Large Alle zusammen:}\\[1cm]
\pause
{\Huge \es{¡Hace sol!}}\\[0.5cm]
\pause
{\Huge \es{¡Hace frío!}}\\[0.5cm]
\pause
{\Huge \es{¡Llueve!}}\\[0.5cm]
\pause
{\Huge \es{¡Nieva!}}
\end{center}
\end{frame}

\begin{frame}{}
\begin{center}
{\Large\bfseries $\rightarrow$ Lernblatt: Kasten 1 ausfüllen}\\[1cm]
{\normalsize Übertrage die Wetter-Vokabeln ins Lernblatt.}
\end{center}
\end{frame}

% ═══════════════════════════════════════════════════════════════
% PHASE 3: ÜBUNG 1 - FOTOS (10 Min)
% ═══════════════════════════════════════════════════════════════

\begin{frame}{}
\phase{Übung 1: Fotos beschreiben}
\end{frame}

\begin{frame}{¿Qué tiempo hace en la foto...?}
\begin{center}
{\Large Buch S. 102-103}\\[1cm]
{\LARGE \es{¿Qué tiempo hace en la foto 1?}}\\[0.5cm]
\pause
{\Large Mögliche Antwort:}\\
{\Large \es{En la foto 1 hace sol y hace calor.}}
\end{center}
\end{frame}

\begin{frame}{Jetzt ihr! Fotos 2-6}
\begin{center}
{\Large Beschreibt das Wetter auf den Fotos 2-6.}\\[1cm]
\begin{tabular}{ll}
Foto 2: & ¿Qué tiempo hace? \\[0.5em]
Foto 3: & ¿Qué tiempo hace? \\[0.5em]
Foto 4: & ¿Qué tiempo hace? \\[0.5em]
Foto 5: & ¿Qué tiempo hace? \\[0.5em]
Foto 6: & ¿Qué tiempo hace? \\
\end{tabular}
\end{center}
\end{frame}

% ═══════════════════════════════════════════════════════════════
% PHASE 4: ÜBUNG 2 - PARTNER-INTERVIEW (15 Min)
% ═══════════════════════════════════════════════════════════════

\begin{frame}{}
\phase{Übung 2: Partner-Interview}
\end{frame}

\begin{frame}{Cuando hace sol, voy...}
\begin{center}
{\Large Die Struktur:}\\[0.5cm]
{\LARGE \es{Cuando} + \textit{Wetter}, + \textit{Aktivität}}\\[1cm]
\pause
{\Large Beispiel:}\\[0.3cm]
{\Large \es{Cuando hace sol, voy al parque con mis amigos.}}\\
{\normalsize \de{Wenn es sonnig ist, gehe ich mit meinen Freunden in den Park.}}
\end{center}
\end{frame}

\begin{frame}{Nützliche Verben und Orte}
\begin{columns}
\column{0.5\textwidth}
\textbf{Verben:}
\begin{itemize}
\item \es{voy al/a la...} \de{ich gehe zu...}
\item \es{me quedo en casa} \de{ich bleibe zu Hause}
\item \es{juego al fútbol} \de{ich spiele Fußball}
\item \es{leo un libro} \de{ich lese ein Buch}
\end{itemize}

\column{0.5\textwidth}
\textbf{Orte:}
\begin{itemize}
\item \es{el parque} \de{der Park}
\item \es{la piscina} \de{das Schwimmbad}
\item \es{el cine} \de{das Kino}
\item \es{la playa} \de{der Strand}
\end{itemize}
\end{columns}
\end{frame}

\begin{frame}{Partner-Interview}
\begin{center}
{\Large\bfseries $\rightarrow$ AB Aufgabe 1}\\[0.8cm]
{\Large Befragt euch gegenseitig:}\\[0.5cm]
\begin{tabular}{l}
\es{¿Qué haces cuando hace sol?} \\[0.5em]
\es{¿Qué haces cuando llueve?} \\[0.5em]
\es{¿Qué haces cuando hace frío?} \\[0.5em]
\es{¿Qué haces cuando hace calor?} \\
\end{tabular}
\\[0.8cm]
{\normalsize Notiert die Antworten eures Partners auf dem Arbeitsblatt.}
\end{center}
\end{frame}

% ═══════════════════════════════════════════════════════════════
% PHASE 5: INPUT ORTE (10 Min)
% ═══════════════════════════════════════════════════════════════

\begin{frame}{}
\phase{Input: Los lugares}
\end{frame}

\begin{frame}{Los lugares \de{– Die Orte}}
\begin{columns}
\column{0.5\textwidth}
\begin{itemize}
\item \es{el parque} \de{der Park}
\item \es{la plaza} \de{der Platz}
\item \es{el cine} \de{das Kino}
\item \es{el teatro} \de{das Theater}
\item \es{la biblioteca} \de{die Bibliothek}
\item \es{la iglesia} \de{die Kirche}
\end{itemize}

\column{0.5\textwidth}
\begin{itemize}
\item \es{el supermercado} \de{der Supermarkt}
\item \es{la panadería} \de{die Bäckerei}
\item \es{la farmacia} \de{die Apotheke}
\item \es{el restaurante} \de{das Restaurant}
\item \es{la cafetería} \de{das Café}
\item \es{el centro comercial} \de{das EKZ}
\end{itemize}
\end{columns}
\end{frame}

\begin{frame}{Tipp: Geschäfte auf -ería}
\begin{center}
{\Large Viele Geschäfte enden auf \es{-ería}:}\\[1cm]
\begin{tabular}{ll}
\es{la panadería} & \de{die Bäckerei} (pan = Brot) \\[0.5em]
\es{la carnicería} & \de{die Metzgerei} (carne = Fleisch) \\[0.5em]
\es{la librería} & \de{die Buchhandlung} (libro = Buch) \\[0.5em]
\es{la pastelería} & \de{die Konditorei} (pastel = Kuchen) \\
\end{tabular}
\end{center}
\end{frame}

\begin{frame}{}
\begin{center}
{\Large\bfseries $\rightarrow$ Lernblatt: Kasten 2 ausfüllen}\\[1cm]
{\normalsize Übertrage die Orte-Vokabeln ins Lernblatt.}
\end{center}
\end{frame}

% ═══════════════════════════════════════════════════════════════
% PHASE 6: ÜBUNG 3 - WETTER + KLEIDUNG (15 Min)
% ═══════════════════════════════════════════════════════════════

\begin{frame}{}
\phase{Übung 3: Wetter + Kleidung}
\end{frame}

\begin{frame}{¿Qué ropa llevas cuando...?}
\begin{center}
{\Large Die Frage:}\\[0.5cm]
{\LARGE \es{¿Qué ropa llevas cuando hace frío?}}\\[0.3cm]
{\normalsize \de{Welche Kleidung trägst du, wenn es kalt ist?}}\\[1cm]
\pause
{\Large Beispielantwort:}\\[0.3cm]
{\Large \es{Cuando hace frío, llevo un abrigo y una bufanda.}}
\end{center}
\end{frame}

\begin{frame}{Kleidung (Wiederholung)}
\begin{columns}
\column{0.5\textwidth}
\textbf{Bei Kälte:}
\begin{itemize}
\item \es{el abrigo} \de{der Mantel}
\item \es{la bufanda} \de{der Schal}
\item \es{los guantes} \de{die Handschuhe}
\item \es{el gorro} \de{die Mütze}
\item \es{las botas} \de{die Stiefel}
\end{itemize}

\column{0.5\textwidth}
\textbf{Bei Wärme/Regen:}
\begin{itemize}
\item \es{la camiseta} \de{das T-Shirt}
\item \es{pantalones cortos} \de{kurze Hosen}
\item \es{las gafas de sol} \de{die Sonnenbrille}
\item \es{el paraguas} \de{der Regenschirm}
\end{itemize}
\end{columns}
\end{frame}

\begin{frame}{Partnerarbeit: Wetter + Kleidung}
\begin{center}
{\Large\bfseries $\rightarrow$ AB Aufgabe 2}\\[0.8cm]
{\Large Schreibt vollständige Sätze:}\\[0.5cm]
\begin{tabular}{l}
a) ¿Qué ropa llevas cuando hace mucho frío? \\[0.5em]
b) ¿Qué ropa llevas cuando hace mucho calor? \\[0.5em]
c) ¿Qué ropa llevas cuando llueve? \\[0.5em]
d) ¿Qué ropa llevas cuando hace buen tiempo? \\
\end{tabular}
\end{center}
\end{frame}

% ═══════════════════════════════════════════════════════════════
% PHASE 7: SPIEL (10 Min)
% ═══════════════════════════════════════════════════════════════

\begin{frame}{}
\phase{Spiel: Wetter-Bingo}
\end{frame}

\begin{frame}{Wetter-Bingo}
\begin{center}
{\Large\bfseries So geht's:}\\[0.8cm]
\begin{enumerate}
\setlength\itemsep{0.5em}
\item Zeichnet ein 3×3 Raster auf ein Blatt.
\item Schreibt 9 verschiedene Wetter-Ausdrücke hinein (Spanisch!).
\item Ich rufe Wetter auf Deutsch – ihr kreuzt das spanische Wort an.
\item Wer zuerst eine Reihe hat, ruft: \es{¡Bingo!}
\end{enumerate}
\end{center}
\end{frame}

\begin{frame}{Wetter-Wörter für Bingo}
\begin{center}
\begin{tabular}{ccc}
hace sol & hace frío & hace calor \\[0.8em]
hace viento & llueve & nieva \\[0.8em]
hay nubes & hace buen tiempo & hace mal tiempo \\
\end{tabular}
\end{center}
\end{frame}

% ═══════════════════════════════════════════════════════════════
% PHASE 8: SICHERUNG (5 Min)
% ═══════════════════════════════════════════════════════════════

\begin{frame}{}
\phase{Sicherung}
\end{frame}

\begin{frame}{Was haben wir heute gelernt?}
\begin{itemize}
\setlength\itemsep{1em}
\item[\textcolor{spanisch}{\checkmark}] \textbf{Wetter:} Hace sol, hace frío, llueve, nieva...
\item[\textcolor{spanisch}{\checkmark}] \textbf{Orte:} el parque, la panadería, el cine...
\item[\textcolor{spanisch}{\checkmark}] \textbf{Struktur:} Cuando hace..., voy a... / llevo...
\end{itemize}

\vspace{1cm}

\begin{block}{Nächste Stunde}
\textbf{Präpositionen:} Wo ist was? \es{enfrente de, al lado de, entre...}
\end{block}
\end{frame}

\begin{frame}{}
\begin{center}
{\Huge\bfseries\textcolor{spanisch}{¡Hasta mañana!}}\\[1cm]
{\Large Bis morgen!}
\end{center}
\end{frame}

\end{document}
